\documentclass[../../main.tex]{subfiles}
\graphicspath{{../../image/}{../../figure/}}
\begin{document}

\section{注记}\label{section:ch06-注记}

关于标量守恒律经典解破裂的更广泛讨论见 Majda [4]。\cref{theorem:ch06-6-1-1} 归功于 Conway [1]。关于激波形成与传播的几何特征的系统研究，见 Izumiya 与 Kossioris [1]，以及 Danilov 与 Mitrovic [3]。\eqref{align:ch06-6-1-1} 到线性输运方程\eqref{align:ch06-6-1-9}的约化是经典的；见 Courant--Hilbert [1,§I.5]。

关于标量守恒律弱解的文献浩如烟海。该研究始于 20 世纪 50 年代，在单空间维数的框架下，受历史导论中已引用的 Hopf [1] 的开创性论文激励。早期工作的参考文献将在\cref{section:ch11-真正非线性失效时}中随相关处给出。


若干空间维数中的第一个存在性证明归功于 Conway 与 Smoller [1]，他们认识到空间 $BV$ 的相关性，并通过 Lax--Friedrichs 差分格式构造了有界变差解。Volpert [1] 后来给出了空间 $BV$ 中解的决定性处理，他显然是第一个在若干空间维数中认识到 L$^{1}$ 压缩性质的人。Kružkov [1] 在 Volpert 工作的基础上，提出了记录在\cref{section:ch06-可容许弱解稳定性}中的可容许弱解刻画，导出 L$^{1}$ 压缩估计，并沿我们在\cref{section:ch06-黏性消失法}中讨论的路线建立了黏性消失法的收敛性。

当流通量关于 $u$ 仅连续时，需要更精细的处理；见 Kruzhkov 与 Panov [1]，Bénilan 与 Kruzhkov [1]，以及 Andreianov、Bénilan 与 Kruzhkov [1]。此外，当流通量不连续时，可容许弱解的概念须澄清并重新定义；见 Panov [8,9,10]，Andreianov、Karlsen 与 Risebro [1]，以及 Audusse 与 Perthame [1]。另一方面，分析可例行推广到非齐次标量平衡律 (3.3.1)，尽管当产生项关于 $u$ 超线性增长时解可能在有限时间 blow up；见 Natalini、Sinestrari 与 Tesei [1]。特别地，具流通量 $G(u,x)=f(u)V(x)$ 的“输运型”非齐次守恒律具有有趣的结构，尤其当 $\operatorname{div}V=0$ 时；见 Caginalp [1] 与 Otto [2]。

为处理流形上的标量守恒律，需要具有几何风味的新思想，因为熵的概念不能以直接方式推广到该框架；见 Amorim、Ben-Artzi 与 LeFloch [1]，Ben-Artzi 与 LeFloch [1]，Ben-Artzi、Falcovitz 与 LeFloch [1]，LeFloch 与 Okutmustur [1]，Dziuk、Kröner 与 Müller [1]，Kröner、Müller 与 Strehlau [1]，以及 Panov [3]。

一般（不必自反）Banach 空间中非线性压缩半群的理论归功于 Crandall 与 Liggett [1]。\cref{section:ch06-压缩半群与周期解}中给出的对标量守恒律的应用取自 Crandall [1]。

可容许解的另一泛函分析刻画见 Portilheiro [1]。

\cref{section:ch06-分层法}中讨论的分层法构造解由 Roždestvenskii [1] 提出，并由 Kuznetsov [1] 与 Douglis [2] 完成。


有一个活跃的研究方向，旨在把双曲守恒律作为更大但更简单的方程组的“松弛”形式来处理，后者可能支配或模拟物理中的松弛现象。进一步讨论与文献见第 XVII 章。\cref{section:ch06-松弛法}的表述遵循 Katsoulakis 与 Tzavaras [1]。

尽管\eqref{align:ch06-6-6-1}是为分析目的人为构造的，但事后可把它解释为描述介观尺度上相互作用粒子系综演化的方程组。见 Katsoulakis 与 Tzavaras [2]，以及 Jin、Katsoulakis 与 Xin [1]。Natalini [2] 中讨论了用松弛格式构造若干维数标量守恒律解的另一途径。

\cref{section:ch06-动力论表述}中所述的动力论表述归功于 Perthame 与 Tadmor [1]，以及 Lions、Perthame 与 Tadmor [2]。

\cref{theorem:ch06-6-7-2} 最初由 Lions、Perthame 与 Tadmor [2] 在 $s<r/(r+2)$ 时建立。改进范围 $s<r/(2r+1)$ 由 Tadmor 与 Tao [1] 借助更尖锐的速度平均估计导出。详细讨论（含推广、应用与广泛文献）见 Perthame [2,3] 的近期专著与综述文章。相关结果见 Giga 与 Miyakawa [1]，Bäcker 与 Dressler [1]，Brenier [1]，James、Peng 与 Perthame [1]，Natalini [2]，Perthame [1]，Perthame 与 Pulvirenti [1]，Hwang [1,2]，Vasseur [2,3,5]，Westdickenberg 与 Noelle [1]，Kissling 与 LeFloch [1]，以及 Dalibard [1]。\cref{theorem:ch06-6-7-2} 中所述诱导正则化效应的机制在一般非线性输运方程（包括经典 Boltzmann 方程，参见 DiPerna 与 Lions [1]）理论中起突出作用。

\cref{section:ch06-动力论表述}中勾勒的可容许解等值集与 Hilbert 空间中压缩半群之间令人惊讶的联系由 Brenier [3] 发现。

另见 Bolley、Brenier 与 Loeper [1]。

关于另一在流体力学极限下松弛到标量守恒律的有趣动力论模型，见 Portilheiro 与 Tzavaras [1]。

构造解还有若干其他方法，最著名的有分数步、谱黏性逼近，或可用于高效计算的各种差分格式。例如见 Bouchut 与 Perthame [1]，Chen、Du 与 Tadmor [1]，Cockburn、Coquel 与 LeFloch [1]，以及 Crandall 与 Majda [1]。数值方面的文献读者应参考 LeVeque [1]，Godlewski 与 Raviart [1,2]，以及 Kröner [1]。

除 L$^{1}$ 与 $BV$ 外，其他函数空间也与理论相关。DeVore 与 Lucier [1] 证明\eqref{align:ch06-6-1-1}的解属于 Besov 空间。

Perthame 与 Westdickenberg [1] 建立了解的总振荡递减性质。

为感受守恒律奇异摄动时解的极限行为，读者可查阅 Botchorishvili、Perthame 与 Vasseur [1]（刚性源效应），Hwang [3]（扩散--色散极限），Aggarwal、Colombo 与 Goatin [1]（非局部效应），以及 Dalibard [2]（齐次化的后果）。

L$^{\infty}$ 解的精细结构，特别地\cref{theorem:ch06-6-8-2}，在 De Lellis、Otto 与 Westdickenberg [1] 中讨论。另见 De Lellis 与 Rivière [1]，De Lellis 与 Golse [1]，以及 Crippa、Otto 与 Westdickenburg [1]。\cref{theorem:ch06-6-8-3} 归功于 Vasseur [4]。另见 Chen 与 Rascle [1]；近期发展见 Panov [5,6] 与 Kwon 与 Vasseur [1]。

\cref{section:ch06-初边值问题}中阐述的用黏性消失法构造初边值问题 $BV$ 解取自 Bardos、Leroux 与 Nédélec [1]。

当 $u$ 与 $\bar{u}$ 仅属于 $L^{\infty}$ 时\cref{theorem:ch06-6-9-3} 的证明见 Otto [1] 与 Málek、Nečas、Rokyta 与 Růžička [1]。该方向近期结果见 Kwon [1]，以及 Coclite、Karlsen 与 Kwon [1]。对流通量在边界为零的初边值问题在 $L^{\infty}$ 中的情形，见 Bürger、Frid 与 Karlsen [1]。流通量不连续的情形由 Carrillo [1] 讨论。L$^{\infty}$ 解经由动力论表述构造，见 Nouri、Omrane 与 Vila [1]，以及 Tidriri [1]。测度值解见 Szepessy [1] 与 Kondo 与 LeFloch [1]。

\eqref{align:ch06-6-1-1}、\eqref{align:ch06-6-1-2} 解的大时间行为在 Conway [1]，Engquist 与 E [1]，Bauman 与 Phillips，以及 Feireisl 与 Petzeltová [1] 中讨论。Chen 与 Frid [1,3,4,6] 建立了在一般守恒律方程组中研究由尺度不变性与紧性诱导的解衰减的框架。特别地，该理论确立了 $u_{0}$ 为周期或形如 $u_{0}(x)=v(|x|^{-1}x)+w(x)$（$w\in L^{1}(\mathbb{R}^{m})$）时\eqref{align:ch06-6-1-1}、\eqref{align:ch06-6-1-2} 解的大时间行为。此框架由 Panov [11,13] 改进并推广到殆周期情形，见 Frid [5] 与 Panov [12]。经由压缩半群通向\cref{theorem:ch06-6-4-9} 的途径取自 Dafermos [35]。与此相关另见 Dafermos 与 Slemrod [1]。

随机效应见 Lions、Perthame 与 Souganidis [2,3]，Gess 与 Souganidis [1,2]，以及这些有趣论文中的文献。

关于 $L^{p}$ 稳定性在可对称化线性组中充要条件\eqref{align:ch06-6-10-4} 的证明归功于 Brenner [1]。Rauch [1] 证明了\cref{theorem:ch06-6-10-1}，Dafermos [22] 证明了\cref{theorem:ch06-6-10-2}。唯一性结果另见 Frid 与 LeFloch [1]。\cref{theorem:ch06-6-10-3} 归功于 Ambrosio 与 De Lellis [1]。另见 Ambrosio、Bouchut 与 De Lellis [1]。最后，Bressan [11,13] 与 De Lellis [1,2] 解释了为何当 $m>1$ 时方程组\eqref{align:ch06-6-10-9}的 Cauchy 问题一般在 $L^{\infty}$ 或 $BV$ 中不适定（\cref{theorem:ch06-6-10-4}）。与之对照，当 $m=1$ 时该方程组的 Cauchy 问题适定并具有有趣理论；见 Temple [2]，Isaacson 与 Temple [1]，Liu 与 Wang [1]，Tveito 与 Winther [1]，Freistühler [7]，以及 Panov [4]。

\end{document}
